% Lecture 10 notes: \input by notes.tex, \subimport'ed by the semester book.
% Exam questions (экзамен): 24 (алгоритм Ленстры факторизации с помощью эллиптической кривой,
% его сложность без доказательства).
\begin{lecture}{10}{Алгоритм Ленстры факторизации с помощью эллиптических кривых}

\section{От алгоритма \texorpdfstring{$p - 1$}{p - 1} к эллиптическим кривым}

Стойкость RSA опирается на трудность разложения модуля $N = pq$ на множители (лекция~5).
Методы факторизации делятся на две группы. \emph{Специальные} методы (пробное деление,
$\rho$-метод и $p - 1$ алгоритм Полларда) работают тем быстрее, чем меньше \emph{наименьший}
простой делитель $p$ числа $N$ (или чем «удобнее» $p$); \emph{общие} методы (квадратичное решето,
лекция~6, решето числового поля) имеют сложность, зависящую только от размера $N$. Алгоритм
Х.~Ленстры~\cite{lenstra-ecm} (1985, опубликован в 1987), который по-английски называют
ECM (\emph{elliptic curve method}), "--- лучший известный специальный метод: он находит
простой делитель $p$ за время, субэкспоненциальное в $\log p$. Поэтому, например, в RSA
простые $p$ и $q$ выбирают одной длины: найти множитель из $40$--$60$ десятичных знаков методом
Ленстры сегодня вполне реально.

\subsection{Напоминание: \texorpdfstring{$p - 1$}{p - 1} алгоритм Полларда}

Пусть $p \mid N$ "--- простой делитель, неизвестный нам. Алгоритм Полларда~\cite{pollard-pm1}
(лекция~5) вычисляет $a_k = a^{k!} \bmod N$ для $k = 2, 3, \dots$ (рекуррентно:
$a_k = a_{k-1}^{\,k} \bmod N$) и проверяет $d_k = \gcd(a_k - 1, N)$.

\begin{proposition}\label{l10:pm1}
Если $p \mid N$, $p \nmid a$ и $(p - 1) \mid k!$, то $p \mid d_k$. Если, кроме того, для
другого простого делителя $q \mid N$ выполнено $a^{k!} \not\equiv 1 \pmod q$, то
$1 < d_k < N$.
\end{proposition}

\begin{proof}
По малой теореме Ферма $a^{p-1} \equiv 1 \pmod p$, а $k! = (p - 1) t$, поэтому
$a^{k!} = (a^{p-1})^t \equiv 1 \pmod p$, т.\,е. $p \mid a_k - 1$ (сравнение $a_k \equiv a^{k!}$
по модулю $N$ влечёт его и по модулю $p$). Второе утверждение: $q \nmid a_k - 1$, значит,
$d_k \ne N$.
\end{proof}

Точнее, достаточно, чтобы $k!$ делилось на порядок $\ord_p(a)$ элемента $a$ в группе $\Znx{p}$,
а он делит $\#\Znx{p} = p - 1$. Алгоритм быстр, если $p - 1$ \emph{гладкое}.

\begin{definition}
Целое число $n \ge 1$ называется \emph{$B$-гладким}, если все его простые делители не
превосходят $B$.
\end{definition}

Если $p - 1$ является $B$-гладким и не делится на высокие степени простых, делитель $p$ будет
найден при $k \le B$, т.\,е. за $\bigO(B \log B)$ умножений по модулю $N$. Но если у $p - 1$
есть большой простой делитель (например, $p = 2p' + 1$ с простым $p'$ "--- так называемое
\emph{безопасное} простое), алгоритм $p - 1$ бесполезен: он требует $k \ge p'$, т.\,е.
работает не быстрее пробного деления.

Корень проблемы в том, что группа $\Znx{p}$ \emph{одна}: её порядок $p - 1$ определяется
числом $p$, и если он неудачный, сделать ничего нельзя. Идея Ленстры: заменить $\Znx{p}$
группой точек эллиптической кривой $E(\FF_p)$. Её порядок тоже близок к $p$, но меняется вместе
с кривой; перебирая кривые, мы рано или поздно получим группу \emph{гладкого} порядка.

\checkyourself{Почему в $p - 1$ алгоритме вычисляют $a^{k!}$, а не $a^{k}$? Чем плох показатель
$\lcm(1, 2, \dots, k)$ по сравнению с $k!$ и чем он лучше?}

\section{Эллиптические кривые над \texorpdfstring{$\ZZ/N\ZZ$}{Z/NZ}}

Всюду далее $N$ "--- составное число, $\gcd(N, 6) = 1$ (делители $2$ и $3$ отделяются пробным
делением), и $N$ не является степенью простого (это проверяется извлечением корней
$\sqrt[k]{N}$, $k \le \log_2 N$). Пусть $p, q, \dots$ "--- простые делители $N$.

Напомним (лекция~8), что над полем $\FF_p$, $p > 3$, эллиптическая кривая
$E\colon y^2 = x^3 + ax + b$ с $4a^3 + 27b^2 \ne 0$ вместе с бесконечно удалённой точкой
$\ECpt$ образует абелеву группу $E(\FF_p)$ с законом сложения
\begin{equation}\label{l10:law}
  \lambda = \begin{cases}
    \dfrac{y_2 - y_1}{x_2 - x_1}, & x_1 \ne x_2,\\[2ex]
    \dfrac{3x_1^2 + a}{2y_1}, & P = Q,\ y_1 \ne 0,
  \end{cases}\qquad
  x_3 = \lambda^2 - x_1 - x_2, \quad y_3 = \lambda(x_1 - x_3) - y_1,
\end{equation}
$P + Q = (x_3, y_3)$, а если $x_1 = x_2$, $y_1 = -y_2$, то $P + Q = \ECpt$.

Над кольцом $\ZZ/N\ZZ$ те же формулы можно применять «как есть»: все операции, кроме деления,
определены, а деление на $d$ "--- это умножение на $d^{-1} \bmod N$, если он существует.

\begin{lemma}\label{l10:dichotomy}
Для $d \in \ZZ$ ровно одно из двух: либо $\gcd(d, N) = 1$ и $d$ обратим по модулю $N$, либо
$g = \gcd(d, N) > 1$, и тогда $g$ "--- делитель $N$, причём $g < N$, если $d \not\equiv 0
\pmod N$. Расширенный алгоритм Евклида за $\bigO(\log^2 N)$ битовых операций различает эти
случаи и в первом из них находит $d^{-1} \bmod N$.
\end{lemma}

\begin{proof}
Если $\gcd(d, N) = 1$, то по соотношению Безу $ud + vN = 1$, т.\,е. $ud \equiv 1 \pmod N$.
Обратно, если $ud \equiv 1 \pmod N$, то всякий общий делитель $d$ и $N$ делит $1$. Если
$d \not\equiv 0 \pmod N$, то $\gcd(d, N) = \gcd(d \bmod N, N) \le d \bmod N < N$. Оценка
сложности "--- из лекции~3.
\end{proof}

\begin{definition}
Пусть $a, b \in \ZZ/N\ZZ$, $\gcd(4a^3 + 27b^2, N) = 1$. \emph{Точкой} кривой
$E\colon y^2 = x^3 + ax + b$ над $\ZZ/N\ZZ$ назовём пару $(x, y) \in (\ZZ/N\ZZ)^2$,
удовлетворяющую уравнению по модулю $N$, или символ $\ECpt$. Множество точек обозначим
$E(\ZZ/N\ZZ)$.\footnote{Это множество не является группой относительно описанной ниже
операции: она определена не всегда. Правильный объект "--- точки проективной кривой над
кольцом, которые образуют группу, изоморфную $\prod_{p} E(\FF_p)$ при бесквадратном $N$;
для алгоритма это не нужно.}
\end{definition}

\begin{definition}[псевдосложение]\label{l10:pseudo}
Для $P, Q \in E(\ZZ/N\ZZ)$ операция $P \oplus Q$ определяется так:
\begin{enumerate}
  \item если $P = \ECpt$ или $Q = \ECpt$, результат "--- другая точка;
  \item если $x_1 \equiv x_2$ и $y_1 \equiv -y_2 \pmod N$, результат "--- $\ECpt$;
  \item иначе $d = 2y_1$ при $P = Q$ и $d = x_2 - x_1$ при $P \ne Q$; если
    $g = \gcd(d, N) > 1$, операция \emph{не выполняется}, и мы говорим об \emph{отказе со
    знаменателем $d$}; при $g = 1$ результат вычисляется по формулам~\eqref{l10:law} по модулю $N$.
\end{enumerate}
\end{definition}

По лемме~\ref{l10:dichotomy} отказ "--- это удача: $\gcd(d, N)$ "--- делитель $N$, и он
нетривиален, если $d \not\equiv 0 \pmod N$.

\taskexample{ecm-add-mod-n}

\subsection{Приведение по модулю простого делителя}

Для простого $p \mid N$ обозначим через $\pi_p\colon E(\ZZ/N\ZZ) \to E(\FF_p)$ приведение
координат по модулю $p$: $\pi_p(x, y) = (x \bmod p, y \bmod p)$, $\pi_p(\ECpt) = \ECpt$.
Корректность: приведение уравнения кривой по модулю $p$ даёт уравнение кривой над $\FF_p$, и
она невырождена, так как $p \nmid 4a^3 + 27b^2$. Отображение $\ZZ/N\ZZ \to \FF_p$,
$t \mapsto t \bmod p$, "--- гомоморфизм колец, поэтому оно сохраняет суммы и произведения, а
обратимый по модулю $N$ элемент $d$ переводит в обратимый, причём $(d^{-1} \bmod N) \bmod p =
(d \bmod p)^{-1}$.

\begin{theorem}[приведение сохраняет операцию]\label{l10:reduction}
Пусть $P, Q \in E(\ZZ/N\ZZ)$ и псевдосложение $P \oplus Q$ выполняется без отказа. Тогда для
любого простого $p \mid N$
\[
  \pi_p(P \oplus Q) = \pi_p(P) + \pi_p(Q) \quad \text{в группе } E(\FF_p).
\]
\end{theorem}

\begin{proof}
Обозначим $\bar P = \pi_p(P)$, $\bar Q = \pi_p(Q)$ и разберём случаи определения~\ref{l10:pseudo}.

\emph{Случай 1} очевиден: $\pi_p(\ECpt) = \ECpt$ "--- нейтральный элемент.

\emph{Случай 2}: $x_1 \equiv x_2$, $y_1 \equiv -y_2 \pmod N$, тем более по модулю $p$; значит,
$\bar Q = -\bar P$ и $\bar P + \bar Q = \ECpt = \pi_p(\ECpt)$.

\emph{Случай 3, $P = Q$}. Знаменатель $d = 2y_1$ обратим по модулю $N$, поэтому
$p \nmid 2y_1$, т.\,е. $\bar y_1 \ne 0$ в $\FF_p$. Значит, в $E(\FF_p)$ сумма $\bar P + \bar P$
вычисляется по формуле удвоения, и её наклон $\bar\lambda = (3\bar x_1^2 + a)(2\bar y_1)^{-1}$
есть приведение $\lambda \bmod p$. Формулы для $x_3, y_3$ полиномиальны по $\lambda, x_i, y_i$,
поэтому их приведение по модулю $p$ совпадает с координатами $2\bar P$.

\emph{Случай 3, $P \ne Q$}. Знаменатель $d = x_2 - x_1$ обратим по модулю $N$, поэтому
$\bar x_1 \ne \bar x_2$: точки $\bar P$, $\bar Q$ различны и не противоположны, их сумма
в $E(\FF_p)$ вычисляется по формуле хорды, наклон которой "--- приведение $\lambda$. Дальше
как в предыдущем случае.
\end{proof}

\begin{figure}[ht]
\centering
\begin{tikzpicture}[x=1cm,y=1cm,font=\small]
  \node[draw,rounded corners,inner sep=6pt] (N) at (0,0)
    {$E(\ZZ/N\ZZ)$: \ $P,\ 2P,\ 3P,\ \dots$ \ (псевдосложение $\oplus$)};
  \node[draw,rounded corners,inner sep=6pt] (p) at (-3.3,-2) {$E(\FF_p)$, \ $\#E(\FF_p) = p + 1 - t_p$};
  \node[draw,rounded corners,inner sep=6pt] (q) at (3.3,-2) {$E(\FF_q)$, \ $\#E(\FF_q) = q + 1 - t_q$};
  \draw[-{Stealth}] (N) -- node[left=4pt] {$\pi_p$} (p);
  \draw[-{Stealth}] (N) -- node[right=4pt] {$\pi_q$} (q);
\end{tikzpicture}
\caption{Вычисления по модулю $N$ «одновременно» выполняются в группах $E(\FF_p)$ и $E(\FF_q)$,
хотя сами эти группы нам неизвестны: известно лишь $N$.}
\label{l10:fig-crt}
\end{figure}

Умножение точки на число $m$ мы будем выполнять удвоением-сложением (лекция~9): каждый шаг
"--- псевдосложение уже вычисленных кратных $iP \oplus jP$. Индукцией по шагам из
теоремы~\ref{l10:reduction} получаем

\begin{corollary}\label{l10:mult}
Если вычисление $mP$ над $\ZZ/N\ZZ$ удвоением-сложением прошло без отказов и дало точку $R$,
то $\pi_p(R) = m\,\pi_p(P)$ в $E(\FF_p)$ для каждого простого $p \mid N$.
\end{corollary}

\begin{theorem}[основная идея алгоритма Ленстры]\label{l10:main}
Пусть $P \in E(\ZZ/N\ZZ)$, $P \ne \ECpt$, и $p, q$ "--- различные простые делители $N$, причём
порядок $\pi_p(P)$ в группе $E(\FF_p)$ делит $m$, а порядок $\pi_q(P)$ в $E(\FF_q)$ не делит $m$.
Тогда вычисление $mP$ над $\ZZ/N\ZZ$ (удвоением-сложением) обязательно заканчивается отказом,
т.\,е. на некотором шаге встретится знаменатель $d$ с $\gcd(d, N) > 1$.
\end{theorem}

\begin{proof}
Допустим, отказа не было и получилась точка $R$. По следствию~\ref{l10:mult}
$\pi_p(R) = m\,\pi_p(P) = \ECpt$. Но приведение аффинной точки $(x, y)$ "--- аффинная точка,
поэтому $R = \ECpt$. Тогда и $m\,\pi_q(P) = \pi_q(R) = \ECpt$, т.\,е. порядок $\pi_q(P)$ делит
$m$ "--- противоречие.
\end{proof}

Что именно происходит в момент отказа? Пусть знаменатель $d$ делится на $p$. Если шаг "---
удвоение, то $\bar y_1 = 0$: точка $\pi_p$ имеет порядок $2$ и её удвоение равно $\ECpt$ в
$E(\FF_p)$. Если шаг "--- сложение, то $\bar x_1 = \bar x_2$, т.\,е. $\pi_p(P) = \pm\pi_p(Q)$:
чаще всего $\pi_p(P) + \pi_p(Q) = \ECpt$, реже (при $\pi_p(P) = \pi_p(Q)$) формула хорды
применена к совпадающим точкам. В обоих случаях \emph{по модулю $p$ «происходит что-то
особенное»}, а по модулю $q$, как правило, нет "--- и $\gcd(d, N)$ отделяет $p$ от $q$.

\taskexample{ecm-double-fail}

\begin{remark}
Отказ с $\gcd(d, N) = N$ возможен: например, $d = x_2 - x_1 \equiv 0 \pmod N$, когда
$mP \equiv \ECpt$ \emph{сразу} по модулю всех простых делителей. Тогда нужно взять другую
кривую (или другую точку). Это аналог ситуации $d_k = N$ в алгоритме $p - 1$, но здесь он
маловероятен: порядки $\pi_p(P)$ и $\pi_q(P)$ "--- порядки элементов разных, никак не связанных
групп.
\end{remark}

\checkyourself{Пусть $x_1 \equiv x_2 \pmod N$, но $y_1 \not\equiv \pm y_2 \pmod N$. Почему
тогда $y_1^2 \equiv y_2^2 \pmod N$ и как в этом случае найти делитель $N$
(вспомните лекцию~7)?}

\checkyourself{Почему в теореме~\ref{l10:main} нельзя утверждать, что $\gcd(d, N) = p$?
Приведите сценарий, в котором отказ происходит на промежуточном шаге из-за делителя $q$.}

\section{Порядок группы \texorpdfstring{$E(\FF_p)$}{E(F\_p)} и выбор кривой}

\begin{theorem}[Хассе; без доказательства, см. лекцию~8]\label{l10:hasse}
Для эллиптической кривой $E$ над $\FF_p$
\[
  \#E(\FF_p) = p + 1 - t, \qquad \abs{t} \le 2\sqrt{p}.
\]
\end{theorem}

Таким образом, $\#E(\FF_p)$ лежит в \emph{интервале Хассе} $[p + 1 - 2\sqrt p,\ p + 1 + 2\sqrt p]$
длины $4\sqrt p$. Какие значения из него реализуются и как часто?

\begin{theorem}[Дойринг; без доказательства]\label{l10:deuring}
Для каждого целого $t$, $\abs{t} \le 2\sqrt p$, существует эллиптическая кривая над $\FF_p$
с $\#E(\FF_p) = p + 1 - t$.
\end{theorem}

Для алгоритма важно не только существование, но и \emph{распределение} порядков. Его оценил
сам Ленстра.

\begin{theorem}[Ленстра~\cite{lenstra-ecm}; без доказательства]\label{l10:distribution}
Существует эффективно вычислимая константа $c > 0$ со следующим свойством. Пусть $p > 3$
"--- простое, $S$ "--- множество целых чисел из интервала $(p + 1 - \sqrt p,\ p + 1 + \sqrt p)$.
Если $a, b \in \FF_p$ выбираются случайно и равновероятно (среди пар с $4a^3 + 27b^2 \ne 0$), то
\[
  \Pr\bigl[\#E_{a,b}(\FF_p) \in S\bigr] \ge c\,\frac{\#S - 2}{\sqrt p \,\log p}.
\]
\end{theorem}

Иными словами, с точностью до множителя $\bigO(\log p)$ порядок случайной кривой распределён в
центральной половине интервала Хассе не хуже равномерного. В частности, если среди чисел этого
интервала доля $B$-гладких равна $\sigma$, то случайная кривая имеет $B$-гладкий порядок с
вероятностью не меньше $c'\sigma/\log p$, и примерно через $\log p/(c'\sigma)$ кривых мы её
встретим. В $p - 1$ алгоритме такой свободы нет: число $p - 1$ одно.

\subsection{Как выбрать кривую и точку}

Чтобы найти точку на кривой $y^2 = x^3 + ax + b$ по модулю $N$, пришлось бы извлекать квадратный
корень по модулю составного $N$, а это так же трудно, как разложить $N$ на множители
(лекция~7). Поэтому делают наоборот: выбирают случайные $a, x_0, y_0 \in \ZZ/N\ZZ$ и
\emph{вычисляют} $b$, при котором точка $P = (x_0, y_0)$ лежит на кривой:
\[
  b = y_0^2 - x_0^3 - a x_0 \bmod N.
\]
Затем проверяют $g = \gcd(4a^3 + 27b^2, N)$: если $1 < g < N$, делитель уже найден (редкая
удача), если $g = N$ "--- выбирают другие параметры. По модулю каждого $p \mid N$ получается
«случайная» кривая со «случайной» точкой.

\checkyourself{Почему нельзя ограничиться кривыми с $a = 0$ или с $b = 0$ (у них порядки
$\#E(\FF_p)$ подчиняются специальным закономерностям)? Подсказка: что будет с распределением
порядков при $p \equiv 2 \pmod 3$ для кривых $y^2 = x^3 + b$?}

\section{Алгоритм}

\begin{algorithm}[Ленстра, первая стадия]\label{l10:alg}
\begin{algorithmic}[1]
\Require составное $N$, $\gcd(N, 6) = 1$, $N$ "--- не степень простого; граница $B$;
  число кривых $C$.
\Ensure нетривиальный делитель $N$ или «неудача».
\For{$i = 1, \dots, C$}
  \State выбрать случайные $a, x_0, y_0 \in \ZZ/N\ZZ$; \ $b \gets y_0^2 - x_0^3 - a x_0 \bmod N$
  \State $g \gets \gcd(4a^3 + 27b^2, N)$
  \If{$1 < g < N$} \Return $g$ \EndIf
  \If{$g = N$} перейти к следующей кривой \EndIf
  \State $Q \gets P = (x_0, y_0)$
  \For{$j = 2, 3, \dots, B$}
    \State $Q \gets j\,Q$ \Comment{удвоение-сложение над $\ZZ/N\ZZ$, теперь $Q = j!\,P$}
    \If{произошёл отказ со знаменателем $d$}
      \State $g \gets \gcd(d, N)$
      \If{$g < N$} \Return $g$ \EndIf
      \State перейти к следующей кривой \Comment{$g = N$}
    \EndIf
  \EndFor
\EndFor
\State \Return «неудача» \Comment{увеличить $B$ или $C$}
\end{algorithmic}
\end{algorithm}

\begin{theorem}[корректность]\label{l10:correct}
Если алгоритм вернул число $g$, то $g$ "--- нетривиальный делитель $N$. Если для некоторой
кривой порядок $\pi_p(P)$ является делителем $B!$, а порядок $\pi_q(P)$ "--- нет, то на этой
кривой алгоритм остановится с отказом (возможно, с $g = N$).
\end{theorem}

\begin{proof}
Первое утверждение "--- лемма~\ref{l10:dichotomy}: возвращаемое $g$ делит $N$ и $1 < g < N$.
Второе "--- теорема~\ref{l10:main} с $m = B!$ (все вычисления $Q \gets jQ$ вместе составляют
одно вычисление $B!\,P$ удвоениями-сложениями). Остановка может произойти и раньше, при
$j < B$.
\end{proof}

Порядок $\pi_p(P)$ делит $\#E(\FF_p)$, поэтому условие заведомо выполнено, если $\#E(\FF_p)$
"--- $B$-гладкое число, не делящееся на большие степени простых. Вместо $j!$ часто используют
$k = \lcm(1, 2, \dots, B) = \prod_{\ell \le B} \ell^{\floor{\log_\ell B}}$: это число меньше
(его логарифм $\sim B$, а $\log B! \sim B \log B$), но делится на любое $B$-гладкое число,
все примарные делители $\ell^e$ которого не превосходят $B$.

\taskexample{ecm-lenstra}

\begin{example}\label{l10:ex-safe}
Пусть $p = 10007 = 2 \cdot 5003 + 1$ и $q = 100043 = 2 \cdot 50021 + 1$ "--- безопасные
простые, $N = pq = 1001130301$. Алгоритм $p - 1$ с основанием $2$ находит делитель $10007$ лишь
при $k = 5003$. Возьмём кривые $y^2 = x^3 + ax + 1$ с точкой $P = (0, 1)$ (здесь $b = 1$
при любом $a$) и вычислим $Q_k = k!\,P$ по модулю $N$. Порядки точек и номер $k$, при котором
произошёл отказ:
\[
\begin{array}{c|l|l|r|c}
  a & \ord \pi_p(P) & \ord \pi_q(P) & k & \text{делитель} \\ \hline
  1 & 61 & 2^2 \cdot 3^2 \cdot 47 \cdot 59 & 59 & q \\
  2 & 2 \cdot 23 \cdot 109 & 2 \cdot 3 \cdot 8369 & 109 & p \\
  3 & 3359 & 7^3 \cdot 97 & 97 & q \\
  4 & 2 \cdot 5^2 \cdot 67 & 3 \cdot 13 \cdot 2579 & 67 & p \\
  5 & 2 \cdot 5 \cdot 1019 & 2^2 \cdot 3 \cdot 7 \cdot 397 & > 300 & \\
  6 & 2 \cdot 37 \cdot 67 & 3^2 \cdot 7 \cdot 1597 & 67 & p \\
  7 & 2 \cdot 11 \cdot 227 & 2^4 \cdot 3109 & 227 & p \\
  8 & 2 \cdot 3^3 \cdot 23 & 7^2 \cdot 17 & 17 & q \\
\end{array}
\]
Уже при $a = 8$ и $k = 17$: $\ord \pi_q(P) = 833 = 7^2 \cdot 17$ делит $17!$, а
$\ord \pi_p(P) = 1242 = 2 \cdot 3^3 \cdot 23$ не делит; на шаге $Q_{17} = 16\,Q_{16} + Q_{16}$
знаменатель $x_2 - x_1$ делится на $q$, и $\gcd(d, N) = 100043$. Вычисление заняло $16$ умножений
точки на $j \le 17$ "--- сравните с $5003$ возведениями в степень в алгоритме $p - 1$.
При $a = 1$ делитель $q$ найден при $k = 59$, хотя порядок по модулю $p$ (равный $61$) ещё
меньше, чем по модулю $q$: важна не величина порядка, а величина его наибольшего простого
делителя.
\end{example}

Все числа примера проверены в Sage (порядки "--- методом \texttt{order()} точек
на кривых над $\FF_p$ и $\FF_q$).

\taskexample{ecm-lenstra-gcd-n}

\checkyourself{В примере~\ref{l10:ex-safe} при $a = 11$ порядок $\pi_p(P)$ равен
$1250 = 2 \cdot 5^4$, т.\,е. делит $20!$, но не $19!$. Тем не менее делитель $p$ находится уже
при $k = 17$, на шаге $16\,Q_{16} + Q_{16}$. Объясните: какое равенство выполняется по модулю
$p$ и почему знаменатель $x_2 - x_1$ делится на $p$?}

\section{Сложность}

\subsection{Стоимость одной кривой}

Вычисление $Q \gets jQ$ "--- это $\bigO(\log j)$ псевдосложений, каждое "--- $\bigO(1)$
умножений по модулю $N$ и одно обращение (расширенный алгоритм Евклида), т.\,е.
$\bigO(\log^2 N)$ битовых операций. На кривую с границей $B$ уходит $\bigO(B \log B)$
псевдосложений (с показателем $\lcm(1, \dots, B)$ "--- $\bigO(B)$), т.\,е.
$\bigO(B \log^2 N)$ битовых операций с точностью до логарифмического множителя.

\subsection{Гладкие числа}

Обозначим через $\psi(x, y)$ количество $y$-гладких чисел, не превосходящих $x$, и положим
\[
  L_x[\alpha, c] = \exp\bigl((c + o(1))\,(\ln x)^{\alpha} (\ln\ln x)^{1 - \alpha}\bigr),
  \qquad 0 \le \alpha \le 1.
\]
При $\alpha = 0$ это $(\ln x)^{c + o(1)}$ (полиномиальная сложность от длины $x$), при $\alpha = 1$
"--- $x^{c + o(1)}$ (экспоненциальная); промежуточные $\alpha$ дают \emph{субэкспоненциальную}
сложность (лекция~6).

\begin{theorem}[Канфилд"--~Эрдёш"--~Померанс; без доказательства]\label{l10:cep}
Если $u = \ln x / \ln y \to \infty$ и $u < (1 - \varepsilon) \ln x/\ln\ln x$, то
$\psi(x, y)/x = u^{-u(1 + o(1))}$.
\end{theorem}

\subsection{Эвристический выбор \texorpdfstring{$B$}{B}}

Пусть $p$ "--- наименьший простой делитель $N$. Предположим (это и есть недоказанная
гипотеза), что порядок $\#E(\FF_p)$ случайной кривой $B$-гладок с той же вероятностью
$u^{-u}$, $u = \ln p / \ln B$, что и случайное число порядка $p$ (теорема~\ref{l10:distribution}
обосновывает это лишь для чисел из интервала Хассе и лишь с точностью до множителя $\log p$, а
гладкость чисел \emph{в коротких интервалах} не доказана). Тогда нужно в среднем $u^u$ кривых по
$\approx B$ операций, и полная работа
\[
  T(B) \approx B \cdot u^u, \qquad \ln T \approx \ln B + u \ln u.
\]
Положим $\ln B = c\sqrt{\ln p \ln\ln p}$. Тогда $u = \frac1c \sqrt{\ln p/\ln\ln p}$,
$\ln u = \frac12 \ln\ln p\,(1 + o(1))$ и
\[
  u \ln u = \frac{1}{2c}\sqrt{\ln p \ln\ln p}\,(1 + o(1)), \qquad
  \ln T = \Bigl(c + \frac{1}{2c} + o(1)\Bigr) \sqrt{\ln p \ln\ln p}.
\]
Функция $c + \frac{1}{2c}$ минимальна при $c = \frac{1}{\sqrt 2}$ и равна $\sqrt 2$. Итак,
оптимальная граница $B = L_p[\frac12, \frac{1}{\sqrt 2}]$, нужно около
$L_p[\frac12, \frac{1}{\sqrt 2}]$ кривых, и

\begin{theorem}[сложность алгоритма Ленстры; без доказательства]\label{l10:complexity}
При естественной гипотезе о гладких числах в интервале Хассе алгоритм Ленстры находит
наименьший простой делитель $p$ числа $N$ в среднем за
\[
  L_p\bigl[\tfrac12, \sqrt 2\bigr] \cdot \bigO(\log^2 N)
  = \exp\Bigl((\sqrt 2 + o(1))\sqrt{\ln p \ln\ln p}\Bigr) \cdot \bigO(\log^2 N)
\]
битовых операций. В худшем случае $p \approx \sqrt N$ это $L_N[\frac12, 1]$ "--- как у
квадратичного решета.
\end{theorem}

Последнее равенство: при $p = \sqrt N$ имеем $\ln p = \frac12 \ln N$, $\ln\ln p \sim \ln\ln N$,
и $\sqrt 2 \sqrt{\frac12 \ln N \ln\ln N} = \sqrt{\ln N \ln\ln N}$.

\begin{remark}
Сложность зависит от $p$, а от $N$ "--- лишь полиномиально, через стоимость арифметики. Это
принципиальное отличие от квадратичного решета и решета числового поля, сложность которых
определяется размером $N$, а не его делителей.
\end{remark}

Грубые оценки без учёта $o(1)$ и полиномиальных множителей ($\log_{10}$ числа операций):
\[
\begin{array}{r|cccc}
  \text{знаков в } p & 20 & 30 & 40 & 60 \\ \hline
  \rho\text{-метод}, \sqrt p & 10 & 15 & 20 & 30 \\
  \text{ECM}, L_p[\frac12, \sqrt 2] & 8{,}2 & 10{,}5 & 12{,}5 & 16{,}0
\end{array}
\]

\checkyourself{Почему оптимальное $B$ и число кривых в эвристике получились одинаковыми
$L_p[\frac12, \frac{1}{\sqrt 2}]$? Как изменится оптимум, если одна кривая стоит $B^2$
операций?}

\section{Вторая стадия и практика}

\paragraph{Вторая стадия.}
После первой стадии с границей $B_1$ имеем $Q = kP$, $k = \lcm(1, \dots, B_1)$. Часто
$\#E(\FF_p) = s\ell$, где $s$ "--- $B_1$-гладкое, а $\ell$ "--- одно простое число из
$(B_1, B_2]$, $B_2 \approx 100 B_1$. Тогда $\ord \pi_p(Q) \mid \ell$, и достаточно для каждого
простого $\ell_i \in (B_1, B_2]$ проверить, не равна ли $\ell_i Q$ точке $\ECpt$ по модулю $p$.
Разности $\delta_i = \ell_{i+1} - \ell_i$ соседних простых малы и чётны, поэтому, заготовив
точки $2Q, 4Q, \dots, DQ$, вычисляют
\[
  \ell_{i+1} Q = \ell_i Q \oplus \delta_i Q
\]
за \emph{одно} псевдосложение на каждое простое (вместо $\bigO(\log \ell)$). Обращений тоже можно
избежать: вместо проверки каждого знаменателя накапливают произведение $\prod (x(\ell_i Q) -
x(R_j)) \bmod N$ и в конце вычисляют один $\gcd$ с $N$. Вторая стадия заметно увеличивает вероятность
успеха одной кривой ценой умеренного увеличения времени её обработки.

\paragraph{Без обращений.}
Обращение по модулю $N$ в десятки раз дороже умножения. Поэтому на практике точки хранят в
проективных координатах $(X : Y : Z)$ (тогда формулы содержат только умножения), а в конце
стадии вычисляют $\gcd(Z, N)$: точка равна $\ECpt$ по модулю $p$ ровно тогда, когда $p \mid Z$.
Ещё быстрее кривые в форме Монтгомери $By^2 = x^3 + Ax^2 + x$ и кривые Эдвардса, для которых
есть формулы, использующие только координаты $X, Z$. Программа GMP-ECM реализует все эти
приёмы.

\paragraph{Где применяется.}
ECM "--- стандартный способ найти «средние» простые делители: множители из $20$--$40$
десятичных знаков находятся за секунды"--~часы, рекордные найденные этим методом множители
содержат более $80$ знаков. Типичная стратегия полного разложения:
пробное деление $\to$ $\rho$-метод и $p - 1$ $\to$ ECM с растущими границами $\to$ квадратичное
решето или решето числового поля для оставшейся части. ECM особенно полезен для чисел вида
$a^n \pm 1$, у которых много делителей разного размера: с его помощью Р.~Брент завершил
разложение чисел Ферма $F_{10}$ и $F_{11}$. Для RSA с простыми одной длины (по $1024$ бита у
модуля $2048$ бит) ECM бесполезен: $L_p[\frac12, \sqrt 2] \approx 10^{42}$ при $p \approx
10^{308}$, тогда как решето числового поля требует около $10^{35}$ операций.

\begin{center}
\renewcommand{\arraystretch}{1.25}
\begin{tabular}{l|l|l|l}
  метод & сложность & от чего зависит & когда применять \\ \hline
  пробное деление & $\bigO(p)$ & $p$ & $p < 10^{6}$--$10^{8}$ \\
  $p - 1$ Полларда & $\approx B \log B$ & гладкость $p - 1$ & если повезёт с $p - 1$ \\
  $\rho$-метод Полларда & $\bigO(\sqrt p)$ & $p$ & $p < 10^{15}$--$10^{20}$ \\
  ECM (Ленстра) & $L_p[\frac12, \sqrt 2]$ & $p$ & $p$ до $40$--$60$ знаков \\
  квадратичное решето & $L_N[\frac12, 1]$ & $N$ & $N$ до $\approx 100$ знаков \\
  решето числового поля & $L_N[\frac13, \sqrt[3]{64/9}]$ & $N$ & $N > 100$ знаков
\end{tabular}
\end{center}

Здесь $p$ "--- наименьший простой делитель $N$; оценки для $p - 1$, ECM и решёт эвристические
(для ECM "--- при гипотезе о гладких числах в интервале Хассе).

\checkyourself{Почему ECM не является угрозой для RSA, а $p - 1$ алгоритм заставляет избегать
простых, у которых $p - 1$ гладкое, хотя ECM строго сильнее? Подсказка: что нужно для успеха
каждого из методов и с какой вероятностью это выполняется для случайного $p$?}

\checkyourself{Как, имея программу ECM, проверить, является ли найденный делитель $g$ простым,
и что делать с $N/g$?}

\checkyourselfsection

\begin{furtherreading}
  \readitem[разд.~6.6]{hps}{алгоритм Ленстры, пример с подробными вычислениями; разд.~3.5 "---
    алгоритм $p - 1$, разд.~3.7 "--- гладкие числа и функция $L$}
  \readitem[разд.~7.4]{crandall-pomerance}{метод эллиптических кривых: проективные координаты,
    кривые Монтгомери, вторая стадия, эвристика сложности}
  \readitem[гл.~VI, \S\,4]{koblitz}{факторизация с помощью эллиптических кривых}
  \readitem[разд.~7.1]{washington}{алгоритм Ленстры, связь с алгоритмом $p - 1$}
  \readitem[разд.~10.3]{cohen}{ECM с практическими рекомендациями по выбору $B_1, B_2$}
  \readitem[гл.~3]{vasilenko}{методы факторизации, в том числе метод Ленстры (на русском)}
  \readitem{lenstra-ecm}{оригинальная статья: алгоритм, распределение порядков, оценка сложности}
  \readitem{mit-18783}{конспекты курса Э.~Сазерленда: гладкие числа и метод эллиптических
    кривых}
\end{furtherreading}
\end{lecture}
